% ═══════════════════════════════════════════════════════════════
% Section 7: Governance & Neuroethics
% ═══════════════════════════════════════════════════════════════

\section{Governance and Neuroethics}
\label{sec:governance}

This section sets out why a security framework for neural devices needs neuroethics, where
neuroethics enters the scoring and the labels, and how the framework relates to policy
instruments. None of it has been examined by an ethicist, clinician or review board.

\subsection{Why Neuroethics}
\label{sec:why-neuroethics}

Neuroethics asks what may rightly be done to and with a human brain. Roskies divided the
field into the ethics of neuroscience, which asks how brain research and technology should
be conducted and used, and the neuroscience of ethics~\cite{roskies2002neuroethics}; this
framework is concerned only with the first.

A security framework needs it because the harms of a compromised neural device are not all
harms to systems or data. Consent to a medical device can be partial and can be withdrawn;
some effects on tissue or neural function cannot be undone; and a recording of brain
activity cannot be reissued. Conventional scoring has no vocabulary for these, and
neuroethics supplies one. Ienca and Andorno~\cite{ienca2017neurorights} proposed four
rights---cognitive liberty, mental privacy, mental integrity and psychological
continuity---and Yuste et al.~\cite{yuste2017four} set out ethical priorities for
neurotechnology and artificial intelligence.

The rights are contested. Bublitz argues against inflating the catalogue of rights, while
describing Ienca and Andorno's four as better grounded in existing
scholarship~\cite{bublitz2022freedom}; Ienca's own later review warns against the same
inflation~\cite{ienca2021neurorights}; and mental privacy and mental integrity lack agreed
operational definitions~\cite{kellmeyer2022neurorights}. QIF therefore uses the four rights
as \emph{labels} for what a technique puts at stake, defines only measurable signal
properties, and does not argue that new law is needed.

Neuroethics also limits what a security framework may claim. The project holds eight
guardrails as data (\texttt{datalake/qif-guardrails.json}), each pairing a caution from the
literature with a limit on a component. Among them: activity in a brain region does not
identify a cognitive process~\cite{poldrack2006neuroimaging}, so TARA catalogues physical
interference, not cognitive content; brain data is given more credibility than it has
earned~\cite{racine2005neuroethics}; present devices do not read
thoughts~\cite{ienca2018brainleaks,wexler2019neurohype}, so every technique separates what
is possible now from what is projected; and neuroscience is dual-use, pursued by national
security agencies~\cite{tennison2012neurorights}, so every threat description is defensive
and paired with a control. The evidence rule, the three scoring states and the labelling of
diagnostic codes as category references (Sections~\ref{sec:tara}--\ref{sec:nic}) apply
these guardrails in practice.

\subsection{How Neuroethics Enters the Score}

Three commitments are fixed in the NISS arithmetic (Equation~\ref{eq:niss}):

\begin{enumerate}
  \item \textbf{Consent is a scored metric at full weight.} A technique that changes
        nothing physically still scores if it overrides consent, although it cannot
        score above 3.0 (Section~\ref{sec:niss-limits}). Of the \LiveNissScored{} live
        scored techniques, 68 are rated Explicit, 49 Implicit, 5 Partial and 7 None on CV.
  \item \textbf{Harm that cannot be undone counts for more.} Reversibility and
        Neuroplasticity carry full weight, expressing non-maleficence as a number, and the
        PINS flag applies to any irreversible technique whatever its total.
  \item \textbf{Reading and disruption are scored separately.} CR and CD carry half weight
        each, so together they weigh as much as one other metric.
\end{enumerate}

\noindent Rights labels are then assigned by rule: High or Critical CR earns mental privacy
and cognitive liberty; High or Critical CD or BI earns mental integrity; further labels come
from a technique's category, its category references and short hand-made lists. In the live
registry 142 of \LiveTechniques{} techniques carry at least one label (mental privacy 111,
mental integrity 102, cognitive liberty 87, psychological continuity 79). The labels do not
feed back into the score.

\subsection{Consent Tiers}

Each technique's governance record names a consent tier:

\begin{table}[H]
\centering
\caption{Consent tiers: snapshot (v\SnapshotRegistryVersion{}) and live counts.}
\label{tab:consent}
\small
\begin{tabularx}{\textwidth}{@{}l X r r@{}}
\toprule
\textbf{Tier} & \textbf{Requirement} & \textbf{v\SnapshotRegistryVersion{}} & \textbf{Live} \\
\midrule
Standard   & Standard clinical consent                         & 22 & 42 \\
Enhanced   & Extended disclosure and monitoring                & 51 & 70 \\
IRB        & Full institutional review board / ethics committee review & 28 & 52 \\
Prohibited & Prohibited outside controlled research            & 2  & 9 \\
---        & Not assigned                                      & 6  & 0 \\
\bottomrule
\end{tabularx}
\end{table}

\subsection{What an Institutional Programme Has, and QIF Lacks}

The guardrails above are self-imposed: no one outside the project checks them. The
neuroethics programme of the NIH BRAIN Initiative is a published example of more. It treats
neuroethics as concurrent with tool development rather than as review after the
fact~\cite{ramos2017brainneuroethics}, and rests on four standing activities: a Neuroethics
Working Group of neuroethics and neuroscience experts that advises the programme; funded
neuroethics research; scanning of the research portfolio for the ethical questions it
raises; and workshops on topics that steer the programme~\cite{bianchi2018brainneuroethics}.
A later account describes this as integrating neuroethics into neuroscience rather than
appending it~\cite{ramos2019brainintegrating}.

QIF has none of the first, second or fourth, and performs the third only as the author
applying the guardrails to the author's own work. This is a gap statement, not a plan; it
keeps the four activities visible as the standard the framework falls short of. Nothing here
implies that the NIH, the BRAIN Initiative or its working group has reviewed, endorsed or has
any connection to QIF.

\subsection{Alignment with Policy Instruments}

UNESCO's Recommendation on the Ethics of Neurotechnology~\cite{unesco2025recommendation},
adopted in November 2025, sets out values, principles and policy action areas.
Table~\ref{tab:unesco} maps its elements to QIF components. The mapping is the author's
self-assessment; it has not been reviewed by UNESCO or any third party, and ``mapped'' means
that a component addresses the element, not that the element is met.

\begin{table}[H]
\centering
\caption{Self-assessed mapping of UNESCO Recommendation elements to QIF components.}
\label{tab:unesco}
\small
\begin{tabularx}{\textwidth}{@{}l l l X@{}}
\toprule
\textbf{UNESCO Element} & \textbf{Type} & \textbf{Status} & \textbf{QIF Component} \\
\midrule
Human rights \& dignity        & Value     & Mapped  & Rights labels, consent tiers \\
Health \& well-being           & Value     & Mapped  & PINS flag, biological impact scoring \\
Diversity                      & Value     & Mapped  & Open-source, multi-stakeholder model \\
Sustainability                 & Value     & Mapped  & Apache 2.0 license, post-trial access \\
Professional integrity         & Value     & Mapped  & Transparency audit trail \\
Proportionality                & Principle & Mapped  & Graduated severity scoring \\
Freedom of thought             & Principle & Mapped  & Cognitive liberty label, consent metric \\
Privacy                        & Principle & Mapped  & Cognitive reconnaissance (CR) / disruption (CD) metrics \\
Protection of children         & Principle & Mapped  & Age-tiered assent framework (governance documents) \\
Consumer protection            & Policy    & Mapped  & Default-deny architecture (governance documents) \\
Enhancement regulation         & Policy    & Partial & Technical infrastructure only \\
Workplace protections          & Policy    & Mapped  & Mental privacy scoring \\
Behavioral influence           & Policy    & Mapped  & Consent violation metric (CV) \\
Health and well-being          & Policy    & Mapped  & NISS scoring, PINS flag \\
Oversight \& governance        & Policy    & Mapped  & Regulatory mapping (below) \\
Access \& equity               & Policy    & Mapped  & Open-source release \\
\bottomrule
\end{tabularx}
\end{table}

\noindent The table lists 16 rows. The v1.5 text described them as 15 of 17 elements
``implemented'' with two partial, which did not match its own table; ``implemented'' is
replaced by ``mapped'' because no row has been tested against the Recommendation. Two
components cited (the age-tiered assent framework and default-deny architecture) are
described in the project's governance documents and are not part of the TARA or NISS data
analysed here.

\subsection{Regulatory Mapping}

The framework relates to existing instruments as follows. None requires TARA or NISS, and
no regulator has reviewed either.

\begin{itemize}
  \item \textbf{FDA Section 524B}~\cite{fda524b}: mandatory cybersecurity documentation
        for connected medical devices, including threat modeling; TARA could inform the
        threat model and NISS could sit beside CVSS in it
  \item \textbf{EU MDR 2017/745}~\cite{eumdr2017}: medical device risk management, for
        which TARA offers a neural threat taxonomy
  \item \textbf{ISO 14971}~\cite{iso14971}: risk management for medical devices; the NIC
        offers a structured association from technical failure to categories of harm
  \item \textbf{HIPAA}~\cite{hipaa1996} and \textbf{GDPR}~\cite{gdpr2016}: health and
        sensitive personal data, under which neural data may fall (GDPR Article 9)
  \item \textbf{Colorado Privacy Act}~\cite{colorado2024neuraldata}: amended in 2024 to
        classify neural data as sensitive data
\end{itemize}
